\documentclass[journal]{IEEEtran}

\usepackage{cite}
\usepackage{amsmath,amssymb}
\usepackage{newtxtext,newtxmath}
\usepackage{graphicx}
\usepackage{booktabs}
\usepackage{array}
\usepackage{url}
\usepackage{textcomp}
\usepackage[hidelinks]{hyperref}

\newcommand{\mn}{\ensuremath{-}}
\newcommand{\pt}{\,pt}

% ---- figure helper: framed placeholder if a figure PDF is missing -------------------------------
\newcommand{\figorbox}[3]{%
  \IfFileExists{#1}{\includegraphics[width=#2]{#1}}{%
    \fbox{\parbox[c][#3][c]{\dimexpr#2-2\fboxsep-2\fboxrule\relax}{\centering\footnotesize
    Figure placeholder\\\texttt{\detokenize{#1}}}}}}

\begin{document}

\title{It Was the Learning Rate: Compute-Matched Baselines Erase the Apparent Gains of a Memristor Crossbar Repair Controller}

\author{Prateek~Yadav\thanks{The author is an independent researcher. E-mail: prateekyadav1656@gmail.com.}}

\maketitle

% =================================================================================================
\begin{abstract}
Memristor crossbars lose cells while deployed, and a finite pool of spare cells is the repair budget.
Prior repair methods remap and retrain once from a known fault map, and they are rarely compared with baselines that have equal compute and a tuned learning rate.
We study in simulation (ResNet-9, CIFAR-10, three fault waves, 15\% spares) what repairs such a chip over its life.
Remapping to a frozen reference can overwrite learned weights, which living spares avoid by construction; a pre-registered tuned-rate control, however, finds no accuracy difference (+0.03 points, 3 seeds), and an earlier +2.48 came from a barely working rate.
A forward-only BatchNorm refit stays at chance once the spare pool is exhausted (76.06 points below our controller, 0 of 8 seeds).
Our controller D, which paces spare use and gates retraining on a canary set, beats equal-compute retraining-only by 3.78 points (5 of 5 seeds).
Its apparent 34.0-point lead over a living-spares rival, however, shrinks to 4.2 points once that rival's learning rate is tuned, to 1.06 once compute is matched, and to $-0.005$ in a pre-registered eight-seed replication.
On final accuracy D is not distinguishable from the matched rival or from a one-shot Liu-style retrain, so we compare these designs by what each costs, and results vary across seeds.
Tuned, compute-matched baselines and a replication are therefore needed before a repair controller's margin means anything.
All results are simulated; Section~\ref{sec:discussion} gives the caveats.
\end{abstract}


\begin{IEEEkeywords}
Memristor crossbar, stuck-at faults, spare cells, fault-aware retraining, in-memory computing, reliability, pre-registration.
\end{IEEEkeywords}

% =================================================================================================
\section{Introduction}
\IEEEPARstart{A}{n} analog resistive-memory (RRAM) chip stores neural-network weights as cell conductances.
Cells fail while the chip is deployed.
A failed cell is stuck at a high or a low conductance and cannot be written again.
A chip designer can keep a pool of spare cells and replace failed cells with spares.
The pool is finite.
Each wave of new faults uses part of it.
Spare cells are therefore a budget that the chip spends over its life.

Two repair tools exist.
Remapping moves a failed weight onto a spare cell.
Retraining adjusts the healthy weights so the network works around the faults.
Prior work uses each tool, often once and from a known fault map \cite{liu17dac,chen17date,song21ittrna}.
A deployed chip needs both, repeatedly, across several fault waves.

Those papers compare against baselines that rarely receive the same compute or a tuned learning rate \cite{xia19tcad,song21ittrna}.
It does not ask what happens when remapping and retraining alternate over a chip's life.
Two questions follow.
Do the two tools interact?
And how much of the apparent gain of a clever repair controller survives a fair comparison?

We find that they interact.
Remap restored a spare toward a reference that the chip captured once, before any retraining, so each remap overwrote the compensation that earlier retraining had learned.
We found this by reading the write paths, then measured it.
A simple fix exists.
A remapped spare takes the cell's current weight and then trains in later bursts.
We call these living spares.
We first measured a gain of 2.48 points for living spares over frozen references, but at an under-tuned burst rate. A pre-registered control at the tuned rate (3 seeds, paired within seed) found +0.03 points, which is noise-level, and a lifetime accuracy 0.31 points lower.
The +2.48 is therefore superseded, and we report the overwrite as a code-level observation, not as an accuracy gain.

Living spares make a second question matter: how fast to spend the pool, and when to retrain.
Our controller D spends a per-wave share of the free spares (pacing).
It ranks cells by saliency, refits BatchNorm statistics, and retrains in a short burst only when a held-out canary set shows a loss.
Importance-ranked remap, BatchNorm refit, plateau stop and canary-triggered retraining each have prior art \cite{liu17dac,jung21date,li19rramedy}.
We do not claim them.
We use D to ask whether this machinery buys accuracy when every rival gets the same compute.

\textit{Methodology and rigor.}
A first version of this work reported a 34-point lead over a living-spares rival.
At least 97\% of that lead came from an untuned learning rate and unequal compute.
We rebuilt the evaluation to remove both.
We registered hypotheses and decision rules before the counted runs.
We gave each rival D's exact number of burst steps.
We tuned each arm's learning rate and report both a test-picked and a canary-picked choice.
We audited the raw run data and fixed simulator errors that had changed earlier numbers.
We ran a pre-registered replication of our own best margin, and two pre-registered reviewer baselines.

\textit{Results in brief.}
D beats equal-compute retraining-only by 3.78 points, in 5 of 5 seeds, significant after Holm correction.
This is a sequential result: seeds were extended from 3 to 5 after a first look.
D does not beat the strongest matched rival.
The 5-seed lead of 1.06 points over C$'$-matched did not replicate on three new seeds, and the eight-seed mean is $-0.005$ (6 of 8 seeds ahead, Holm $p=1.0$).
That is not a tie we can prove: D is ahead on 6 of 8 seeds by 0.44 to 1.52 points and behind by about 3 points on the 2 seeds where its gate skipped a burst, and we ran no equivalence test.
A one-shot Liu-style retrain at D's compute is not distinguishable from D on final accuracy (Holm $p=0.66$).
It does not match D on service through life (Section~\ref{sec:e1e2}).
A forward-only BatchNorm repair stays at chance in the tested regime, where the spare pool is exhausted (Section~\ref{sec:e1e2}).
On two of the eight seeds, D's canary gate drifted and skipped a needed burst (Section~\ref{sec:gate}).
We also report settings with no edge or a deficit: VGG6 (mixed), many small waves with a gate fixed to the best-so-far canary (mixed), and calibrated drift (D trails B-matched by 0.65 to 3.54 points; negative with 5\% spares; 3 seeds).

\textit{Contributions.}
\begin{enumerate}
\item We identify, by reading the write paths, that a frozen-reference remap restores the original value over retrained weights (Section~\ref{sec:method}), and we report that a pre-registered tuned-rate control found no measurable accuracy effect of the living-spares fix (+0.03 points, 3 seeds; Section~\ref{sec:overwrite}); our design cannot test the remap-after-retraining case.
\item We show that a BatchNorm refit does not rescue a chip whose spare pool is exhausted (the tested regime): it stays at chance on 8 of 8 seeds (Section~\ref{sec:e1e2}).
\item We give a compute-matched, tuned, pre-registered evaluation that cut a 44.7-point (retraining-only) and a 34.0-point (living-spares rival) apparent lead to 3.78 points (5 seeds, significant, sequential design) and to zero (8 seeds, no equivalence test), respectively, and to no detectable difference from a one-shot Liu-style retrain (Sections~\ref{sec:tuning}, \ref{sec:replication} and~\ref{sec:e1e2}). A matched comparison shows that several designs reach about the same final accuracy, and Table~\ref{tab:tradeoffs} gives what each costs in spares, canary data, tuning and service through life.
\item We report the controller's limits: pacing and gating show no accuracy gain at matched compute, the gate has a drift failure mode, and D has no edge or trails on calibrated drift and with many small waves (Sections~\ref{sec:gate} and~\ref{sec:stress}). We also give a closed-loop test of D on a probe-estimated fault map with its limits (Section~\ref{sec:closedloop}), and a condensed catalogue of simulator pitfalls that we found and fixed (Section~\ref{sec:fidelity}).
\end{enumerate}

The rest of the paper is organized as follows.
Section~\ref{sec:background} reviews related work, Section~\ref{sec:method} describes the method, Section~\ref{sec:methodology} the experimental setup, Section~\ref{sec:results} the results, Section~\ref{sec:discussion} limitations, and Section~\ref{sec:conclusion} concludes.
Figure~\ref{fig:overview} gives the system and the story in one view.

\begin{figure*}[t]
\centering
\figorbox{figs/fig1_overview.pdf}{\textwidth}{2.6in}
\caption{System and story. A chip loses cells in waves; a finite spare pool is the repair budget. Remap to a frozen reference restores the original value over what retraining learned. A living reference keeps it (no accuracy gain was measured at the tuned rate). Controller D spends the pool at a paced rate and retrains in short, canary-gated bursts. This is a diagram; it shows no measured data.}
\label{fig:overview}
\end{figure*}

% =================================================================================================
\section{Background and Related Work}
\label{sec:background}

\subsection{Setting and terms}
We use one term for one thing.
\emph{Spare cell}: a reserve cell that can replace a failed cell.
\emph{Wave}: one fault event.
Faults are permanent: a stuck cell stays at its high-resistance (HRS) or low-resistance (LRS) conductance.
\emph{Burst}: a short retraining run between waves.
\emph{Canary set}: 1,024 held-out training images that serve as the only health signal.
\emph{Arm}: one policy in a comparison.
\emph{Deployment regime}: we train once to a checkpoint, then run inference through $K$ fault waves.
Retraining happens only inside policy bursts.
Labelled data for bursts is assumed to be available on the device or at the edge.

\subsection{Related work}
Table~\ref{tab:positioning} positions the closest work.
We read each paper listed there.
Our claims are limited by the coverage of that reading and of our search.

\emph{Retrain once, from a known map.}
Chen et al.\ and Song et al.\ map weights to a crossbar with redundant rows, then retrain \cite{chen17date,song21ittrna}.
Jin et al.\ rank weights by gradient significance and retrain after mapping \cite{jin20date}.
Liu et al.\ retrain with masks, then remap only the most significant defects onto a capped pool of redundant columns \cite{liu17dac}.
All four work once and from a static map.
The conflict between remapping and retrained weights cannot arise in that order.
Importance-ranked remap is Liu's contribution, not ours.

\emph{Lifetime loops.}
Xia et al.\ alternate on-line fault detection with fault-tolerant training; the journal version adds a threshold on the share of newly detected faults \cite{xia17dac,xia19tcad}.
Tung et al.\ remap whole tasks every epoch from fault-density tests during training \cite{tung23date}.
RRAMedy detects a health change on a deployed chip, diagnoses it, and retrains \cite{li19rramedy,li22tc}.
It rejects spares as too costly.
These papers publish the lifetime detect-then-retrain loop.
We do not claim it.
RRAMedy's mitigation policy, run under an oracle map, is our arm B (Table~\ref{tab:arms}).

\emph{Spares and repair budgets.}
Xia et al.\ describe a post-fabrication, reconfigurable spare pool that is allocated greedily to the worst columns \cite{xia18jetcas}.
Error-correcting pointers allocate greedily as cells fail \cite{schechter10ecp}.
Khaleghi and Rao choose \emph{which} shared spare to consume, to keep later repairs possible \cite{khaleghi18}.
Hamun defers retirement of worn cells until the estimated accuracy loss passes a limit \cite{sabri25hamun}.
ReSAFT sizes redundancy at design time \cite{dorostkar26resaft}.
None of these sets \emph{how many} spares to spend per wave inside a retraining controller.
Greedy spending is our baseline for that question.

\emph{Retraining-free and compensation schemes.}
Jung et al.\ refit BatchNorm statistics without gradient steps \cite{jung21date}.
Fault-Free re-encodes weights at compile time and compensates the largest errors \cite{shin21dac}.
SEC-CiM and SEC-CiM+ use a digital error memory under a fixed budget; the budget is used first come, first served \cite{bommana24itc,bommana26todaes}.
SEC-CiM+ avoids the overwrite conflict by not retraining.
DEAR refreshes a stored error after each weight update during training \cite{bommana25dear}.
Maharjan and Lu re-clamp stuck weights after each optimizer step \cite{maharjan26isqed}.
These are partial precedents for parts of our fix: a freeze of stuck weights, and a refresh of a stored reference.
We claim only their combination with a finite pool, trained spares, and a deployed chip (Section~\ref{sec:method}).

\emph{Map-free repair and fair comparison.}
STAR repairs a network without any fault localization, using Equilibrium Propagation after a one-shot fault injection \cite{khan26star}.
Its learning rule does not port to our backpropagation network.
Our arm B-matched is the unassisted baseline of that class, not STAR itself.
Xia et al.\ report a 1.2x iteration count for their method, and Song et al.\ match iterations for one baseline \cite{xia19tcad,song21ittrna}.
We give every rival the controller's step count and tune each arm's learning rate.

\begin{table*}[t]
\centering
\caption{Closest prior work and what this paper does with it. ``Not claimed'' means we use the element as published.}
\label{tab:positioning}
\small
\begin{tabular}{p{2.6cm}p{5.6cm}p{7.6cm}}
\toprule
Work & What it does & Relation to this paper \\
\midrule
Liu DAC'17 \cite{liu17dac} & One-shot masked retraining, then remap of top defects onto a capped pool; oracle map & Importance-ranked remap: not claimed. Our E1 arm adapts it to waves. \\
Chen DATE'17, ITT-RNA \cite{chen17date,song21ittrna} & Map first, retrain once; trained redundant cell in a differential pair; static & Static, offline. Our setting adds waves and a pool drawn at fault time. \\
Jin DATE'20 \cite{jin20date} & Gradient-significance ranking, masked retraining & Ranking: not claimed. \\
Jung DATE'21 \cite{jung21date} & BatchNorm refit, no gradient step & BN refit: not claimed. Our E2 arm tests it alone. \\
Xia DAC'17, TCAD'19 \cite{xia17dac,xia19tcad} & Periodic detection and tolerant training; fault-share gate & Lifetime loop: not claimed. No spares, no pacing. \\
Xia JETCAS'18 \cite{xia18jetcas} & Reconfigurable spare pool, greedy allocation & Nearest spare-pool precedent. Sets where, not when. \\
Tung DATE'23 \cite{tung23date} & Per-epoch whole-crossbar task remap during training & Online remap; no spares, no retraining bursts. \\
RRAMedy ICCD'19, TC'22 \cite{li19rramedy,li22tc} & Alarm, diagnosis, retraining of a deployed chip; no spares & Generic gate and loop: not claimed. Arm B stands for it. \\
SEC-CiM, SEC-CiM+ \cite{bommana24itc,bommana26todaes} & Digital error compensation under a fixed budget; no retraining & Budget without pacing. Discussion only. \\
DEAR, Maharjan \cite{bommana25dear,maharjan26isqed} & Stored-error refresh; stuck-weight re-clamp & Partial precedents for our fix. \\
Hamun, ECP, Khaleghi, ReSAFT \cite{sabri25hamun,schechter10ecp,khaleghi18,dorostkar26resaft} & Deferred retirement; greedy pointers; which spare; design-time sizing & Greedy spending is our baseline. \\
STAR \cite{khan26star} & Map-free repair with Equilibrium Propagation & B-matched is the same class; not a run of STAR. \\
\bottomrule
\end{tabular}
\end{table*}

% =================================================================================================
\section{Method}
\label{sec:method}

\subsection{The overwrite problem}
A remap replaces a stuck cell by a spare.
In a frozen-reference design, the pool stores a reference conductance when it is built.
Every later remap restores the spare toward that reference.
Between waves, retraining changes the weights of healthy cells so that they compensate for faults.
A later remap then writes the old reference over a cell that retraining had adjusted.
The compensation can be lost, and nothing reports it.
This is a code-level observation; Section~\ref{sec:overwrite} shows that we could not measure an accuracy cost at the tuned rate.

\subsection{Living spares}
A living spare takes the cell's current weight when the remap happens.
It is not pinned to a reference, and it trains in later bursts.
A second change keeps stuck weights fixed during retraining.
After each optimizer step, the value of every stuck cell is restored.
Otherwise gradient, momentum and weight decay drift the stuck weights, and a later remap copies that drift into a spare.
A spare that is remapped can still fail again.
It counts against the pool.
Freezing stuck weights and refreshing references both have partial precedents (Section~\ref{sec:background}).
Our contribution is the combination for a finite pool on a deployed chip with several waves.

\subsection{Controller D}
\label{sec:controllerd}
D acts after each wave.
It uses only quantities that a deployed chip can observe: the fault map (oracle or estimated, Section~\ref{sec:closedloop}) and the canary set.
\begin{enumerate}
\item \textbf{Paced remap.}
At wave $k$ of $N$ expected waves, each layer may spend a share $\rho_k = 1/\max(1, N-k)$ of its free spares.
The last expected wave may spend all that remain.
Cells are ranked by a Fisher saliency score.
\item \textbf{BatchNorm refit.}
D refits BatchNorm statistics in a forward-only pass.
\item \textbf{Canary gate.}
D reads the canary set.
If the drop against a reference is below a margin, D skips retraining.
The reference is the canary reading after the previous wave's repair (Section~\ref{sec:gate}).
\item \textbf{Burst.}
Otherwise D retrains in a burst with fresh SGD and a cosine learning-rate decay.
The burst stops at a plateau of the canary signal or at a step cap.
\end{enumerate}
D needs the expected number of waves $N$.
A wrong belief costs about 0.9 to 1.0 point at the tuned rate (3 seeds, two priors files).
D used 1,375 to 1,500 burst steps per seed on the first five seeds, and 700 to 1,500 over all eight.

\subsection{Arms}
Table~\ref{tab:arms} defines every arm.
B-matched and C$'$-matched receive the same number of burst steps as D on the same seed.
That equal compute is the main control of this study.
E1 and E2 are the two reviewer baselines that we registered.

\begin{table*}[t]
\centering
\caption{Arm definitions. All arms share the same checkpoint and fault maps within a seed (hash-checked). B-matched and C$'$-matched use D's burst steps for the seed.}
\label{tab:arms}
\small
\begin{tabular}{p{2.1cm}p{9.2cm}p{4.2cm}}
\toprule
Arm & Definition & Role \\
\midrule
R0 / R1 & No mitigation / remap only, random order, no retraining & Static baselines \\
B & One fixed burst after every wave; no remap, no spares & Retraining only (RRAMedy-class policy under an oracle map) \\
B-matched & B with D's total burst steps for the seed & Equal-compute retraining-only \\
C & Remap; burst only when a wave overflows the pool; frozen reference & Overwrite bug \\
C$'$ & C with living spares & Spares without a controller \\
C$'$-matched & C$'$ with D's burst steps for the seed & Strongest fair rival \\
D & Controller of Section~\ref{sec:controllerd} & This paper \\
E1 & One Liu-style burst at the last wave; greedy saliency-ordered remap each wave; no pacing, gate or BN refit; burst length equals D's steps; rate 0.5 not re-swept & Closest prior art, adapted to waves \\
E2 & Random-order remap each wave, then BatchNorm refit; zero optimizer steps & What a forward-only repair recovers \\
\bottomrule
\end{tabular}
\end{table*}

% =================================================================================================
\section{Experimental Setup}
\label{sec:methodology}

\subsection{Setup}
The network is a ResNet-9 on CIFAR-10.
We run it in the CrossSim simulator (version 3.2.1) with an 8-bit DAC and a calibrated 8-bit ADC.
The main runs have no read noise, no programming error and no drift error.
Each seed trains one checkpoint on 50,000 images.
The fault-free deployed accuracy is 88.92 percent (5 seeds).
The default test has three waves of 10\% random stuck-at cells.
The cumulative stuck fraction is 27.1\%.
The pool holds 15\% spares, one per cell, without area or routing cost.
The fault map is an oracle in the main runs.
Accuracy is the accuracy after the last wave's burst on the first 5,000 test images (sampling noise about $\pm 0.5$ points).
The main test uses seeds 42, 101, 102, 103 and 104, and E3 adds seeds 105, 106 and 107.
B and B-matched were not run on seeds 105 to 107.
The same seed gives the same checkpoint and the same fault maps to every arm; we check the fault-map hash.
Rate and gate settings are in the logs of each run (\texttt{effective\_params}), not in the stored configuration.

\subsection{Statistics}
We pair arms by seed and use a two-sided paired $t$ test as the primary test.
Holm correction runs inside families fixed in advance: a headline family H of 15 comparisons, a limitation family N of 5, and an exploratory family S of 64.
The E1 and E2 comparisons form a family E of 4, and the lifetime (area under the per-wave curve) comparisons form a family of 28.
Rows in family S are exploratory: we mark them ``significant (exploratory family)''.
We call a result \emph{significant} only if the Holm-adjusted $p$ is below 0.05 and the one-sided sign test is below 0.05.
That needs $n \ge 5$ and all seeds in one direction.
We call a result \emph{consistent} if all seeds agree but it is not significant, \emph{mixed} if seeds disagree, and \emph{noise-level} if the mean is under twice the noise of the mean.
Ten identical reruns differ by an RMS of 0.46 points; the noise of a mean of five seeds is about 0.20.
With $n=3$, no result can be significant.

\subsection{Methodology and rigor}
\label{sec:honest}
Table~\ref{tab:ledger} lists each control and its limit.
We kept four habits.

\emph{We registered before we ran.}
The pre-registration of E1 to E5 was committed before any counted run of them.
It states hypotheses, direction, decision rules and the families for Holm correction.
For E3 it says what counts as a replication (Section~\ref{sec:replication}).
Seeds 103 and 104 were added after a first look at seeds 42, 101 and 102.
Seeds 105 to 107 (E3) were added after the five-seed analysis, so the +1.06 point result on five seeds is an interim look.
E3 (seeds 105 to 107) was registered in commit 018276a before it ran.

\emph{We matched compute and tuned every arm.}
Each rival gets D's burst steps.
Each arm's burst learning rate is swept.
The learning rate was the largest hidden factor we found (Section~\ref{sec:tuning}).

\emph{We audited the raw data.}
An independent pass re-derived the numbers from the stored event logs.
It checked 683 complete deployment lifetimes: 562 passed, 121 raised benign warnings, and none failed.
It also corrected 53 wording errors in our own earlier report.
We list the remaining provenance limits in the table.

\emph{We report failures.}
E3 did not replicate our best margin, and E1 could not be separated from D on final accuracy.
A closed-loop extension was killed at 4\% noise.
We found a failure mode of our own controller's gate after the verdicts were fixed, and we label that diagnosis post-hoc.
A pre-registered variant of D, called D2 (D with an adaptive trigger that stops bursts earlier, about 1,220 against 1,470 steps), failed its main-test criterion, and D stays the controller.
D2 passed its calibrated-scenario criterion in 4 of 4 scenarios.
Section~\ref{sec:results} reports each outcome as it came.

\begin{table*}[t]
\centering
\caption{Methodology ledger. Each row is a control we used, what it protects against, and its known limit.}
\label{tab:ledger}
\small
\begin{tabular}{p{3.0cm}p{6.6cm}p{6.1cm}}
\toprule
Control & What we did & Known limit \\
\midrule
Pre-registration & E1 to E5 (including E3) registered at commit 018276a (2026-10-14) before any counted run; decision rules fixed in advance
& Earlier experiments (first-wave spectral work) were registered after nine runs; see companion paper (in preparation) \\
Sequential design & Seeds 103 and 104 added after a first look at seeds 42, 101 and 102; E3 adds seeds 105 to 107
& The eight-seed test is a second look; its $p$ value is optimistic. D versus B-matched (5 seeds) was also extended after a first look. We print new seeds alone and pooled; B-matched has five seeds \\
Equal compute & B-matched and C$'$-matched get D's exact steps ($\pm 1$) per seed
& Compute, not write cycles; endurance is not costed \\
Learning-rate tuning & Per-arm rate sweeps; D keeps the rate it was tuned at
& ResNet-9 rates were picked on seed-42 \emph{test} accuracy; we print both picks (Table~\ref{tab:lr}). D has more settings (gate margin, plateau rule, priors, rate) than the matched rivals (rate only) \\
Raw-data audit & Numbers re-derived from event logs; fault-map hashes equal across arms within a seed
& No run records its code commit (0 of 683); one snapshot commit, 32aa9ae, holds the code the runs used \\
Fault-map pairing & Same seed, same hash across arms
& The same seed gives different maps on different GPU models (9 of 9 multi-GPU buckets differ); seeds 103 and 104 ran on another GPU than 42 to 102 \\
Noise floor & Rerun RMS 0.46 points; largest single difference 1.22
& A single-seed gap below 1.22 points is not interpretable \\
Failed replications and kills reported & E3 failure, E1 not distinguishable, closed-loop 4\% extension, D2, many-small-wave and S4 losses, gate drift
& The gate-drift diagnosis is post-hoc \\
\bottomrule
\end{tabular}
\end{table*}

\subsection{Simulation fidelity}
\label{sec:fidelity}
A simulator that gives the wrong answer for the right reasons is dangerous.
Our simulator needed many fixes before its numbers meant anything.
We keep the catalogue as a contribution, because other groups use similar tools.
We report only entries with a stored measurement in Table~\ref{tab:bugs}.
We found further errors in our development log, such as a training step that never reached the analog core and a backend whose forward pass ignored conductance.
We do not report effect sizes for those, because a stored probe is still missing.

Three lessons stand out.
First, converter details dominate accuracy.
A fixed input-clipping range in the DAC collapsed the analog network at the float-to-analog handover.
Fixing it moved one run from 45.98 to 89.79 percent and removed handover drops of 58 to 75 points.
That comparison is unpaired (one seed against four), so we use it as a scale, not as an effect size.
An earlier ``model capacity'' conclusion was an artifact of this defect.
Second, the repair logic had its own bugs.
Two remap bugs left a third of the spare pool unused and blocked remap at the last wave.
A test suite of 207 tests passed while a mitigation bug was present, because the tests ran each writer alone and never together with an optimizer step.
Third, comparisons need hygiene.
The same seed gives different fault maps on different GPU models.
Best-of-$N$ maxima over 156 evaluations inflate a summary table.
A device comparison can be vacuous: five of five arms of an STT-MRAM run were identical to the default device.

\begin{table*}[t]
\centering
\caption{Simulator pitfalls with a stored measurement. Each was found by us and fixed or documented. ``Effect'' is the measured change in a reported number.}
\label{tab:bugs}
\small
\begin{tabular}{p{3.4cm}p{5.7cm}p{6.6cm}}
\toprule
Pitfall & Symptom & Measured effect \\
\midrule
DAC input clipping & Analog phase collapses at handover; bit-depth sweeps saturate & One run 45.98 $\rightarrow$ 89.79 percent (unpaired, one seed vs four); handover drop 58--75 points $\rightarrow$ under 1
\\
Capacity misdiagnosis & A ladder of net sizes read as a capacity limit & Conclusion invalidated by the DAC fix
\\
Regime saturation & Density 0.1 every 8 batches & $1-0.9^{118}$ = 100\% stuck; every arm at chance
\\
Remap selection stall & Remapped cells stayed in the stuck mask & Pre-fix runs used 656,868 of 985,293 spares (66.7\%, with the gate bug); post-fix runs use all
\\
Synthetic health gate & Remap gated on a drifting stress metric & Remap could never run at wave 3
\\
Stale reference (the overwrite) & Remap restores a reference captured once & Living minus frozen $+2.48$ points at the old rate (superseded); $+0.03$ at the tuned rate, noise-level
\\
Dispatch on healthy cells & Soft mitigation acted on non-stuck cells & About 26 million cells touched per run; a claimed $+2.56$ gain was confounded
\\
Baseline under-tuning & Burst rate at the edge of the grid & Retraining-only $+44.7$ points when tuned, 3 of 3 seeds
\\
Fault-map pairing & Classic arms use the global RNG & 14 of 223 report pairs unpaired
\\
Reporting & Best-of-$N$ maxima; duplicate directory counted as a seed & Maximum over 156 evaluations is optimistic by construction
\\
Test-suite blind spot & Writers tested without an optimizer step & 207 tests passed with a mitigation bug present
\\
Vacuous device comparison & STT-MRAM equals default device & 5 of 5 arms identical; no device conclusion
\\
\bottomrule
\end{tabular}
\end{table*}

% =================================================================================================
\section{Results}
\label{sec:results}

\subsection{Main contrasts}
\label{sec:main}
Figure~\ref{fig:main} and Table~\ref{tab:main} give the main test.
The test has eight seeds for D, C$'$-matched, E1, E2 and R1.
B and B-matched have five seeds, because E3 did not run them on the three new seeds.
Three groups of arms appear.
Remap-only (R1) and the forward-only BatchNorm repair (E2) end at chance.
B-matched, the equal-compute retraining-only arm, ends near 85 percent.
D, C$'$-matched and the one-shot Liu-style arm (E1-last) end within 0.6 points of each other on the mean and are not distinguishable on 8 seeds (no equivalence test).

D beats B-matched by 3.78 points (per seed +4.14, +3.52, +3.68, +4.46, +3.10; 5 seeds).
The 95\% interval is +3.1 to +4.4.
The Holm-adjusted $p$ is 0.0014, so the result is significant.
This is a sequential design: we extended the seeds from 3 to 5 after a first look, so the label is optimistic.
Without seed 42, the seed where rates were picked, the gap is 3.69 in 4 of 4 seeds.
On these five seeds D ends at 88.50 percent, B-matched at 84.72, and the fault-free chip at 88.92.
This is a five-seed result: B-matched was not run on seeds 105 to 107.
Remap alone (R1) ends near chance: D beats it by 78.77 points (5 of 5 seeds, significant, exploratory family S).
Retraining is therefore needed at 15\% spares.
Spares add to retraining: C$'$-matched beats B-matched by 2.72 points (5 of 5 seeds, Holm 0.004, significant in the exploratory family S).
That contrast compares spares and remap with retraining only, so it does not isolate living spares.

\begin{figure}[t]
\centering
\figorbox{figs/fig3_main.pdf}{\columnwidth}{2.1in}
\caption{Final test accuracy per arm after the last wave (ResNet-9, CIFAR-10, three waves of 10\%, 15\% spares, oracle map). Each marker is one seed; the black bar is the mean. Eight seeds (42, 101--107) for D, C$'$-matched, E1, E2 and R1; five seeds (42, 101--104) for B and B-matched; three for R0. Means: D 86.64, C$'$-matched 86.64, E1-last 86.08, B-matched 84.72, E2 10.58. Notice that D and C$'$-matched have the same mean, and that the two low D seeds (105, 106) are the ones that pull D down. The axis is broken between 30 and 79.}
\label{fig:main}
\end{figure}

\begin{table*}[t]
\centering
\caption{Main paired contrasts, ResNet-9, CIFAR-10, 15\% spares (final accuracy, points, unless marked lifetime). Per-seed values follow seed order 42, 101, 102, 103, 104 (and 105, 106, 107 where $n=8$). Fam.\ is the Holm family: H (headline, 15), E (reviewer baselines, 4), S (exploratory, 64). Holm $p$ is within that family. Rows in family S are exploratory, so ``significant'' there is not a headline claim. D $-$ B-matched is a sequential design (seeds extended from 3 to 5 after a first look). The two lifetime rows are exploratory and not pre-registered; their $p$ values are unadjusted. Status follows Section~\ref{sec:methodology}. Wins count seeds where the first arm is ahead.}
\label{tab:main}
\small
\begin{tabular}{p{3.0cm}p{0.4cm}p{0.95cm}p{4.6cm}p{1.35cm}p{0.55cm}p{1.15cm}p{2.3cm}}
\toprule
Contrast (A minus B) & $n$ & Mean & Per seed & Wins & Fam. & Holm $p$ & Status \\
\midrule
D $-$ B-matched & 5 & +3.78 & +4.14, +3.52, +3.68, +4.46, +3.10 & 5/5 & H & 0.0014 & significant (sequential design) \\
D $-$ C$'$-matched, first five seeds & 5 & +1.06 & +1.28, +0.44, +1.18, +1.52, +0.88 & 5/5 & H & 0.056 & consistent (interim look) \\
D $-$ C$'$-matched, new seeds & 3 & $-$1.78 & $-$3.12, $-$3.32, +1.10 & 1/3 & --- & --- & mixed \\
D $-$ C$'$-matched, all seeds (E3) & 8 & $-$0.005 & +1.28, +0.44, +1.18, +1.52, +0.88, $-$3.12, $-$3.32, +1.10 & 6/8 & H & 1.000 & mixed; failure; not distinguishable (no equivalence test) \\
E1-last $-$ D & 8 & $-$0.55 & $-$1.66, $-$1.90, $-$2.06, $-$1.86, $-$0.88, +2.02, +2.82, $-$0.90 & 2/8 (E1 ahead) & E & 0.66 & mixed; not distinguishable (no equivalence test) \\
E2 $-$ D & 8 & $-$76.06 & $-$78.20, $-$79.98, $-$77.20, $-$78.42, $-$77.86, $-$70.36, $-$67.12, $-$79.34 & 0/8 & E & $<$0.001 & negative \\
E2 $-$ R1 & 8 & +0.67 & +0.40, $-$2.18, +0.94, $-$0.02, +3.04, +0.36, +3.38, $-$0.54 & 5/8 & E & 0.66 & mixed \\
E1-last $-$ R1 & 8 & +76.18 & --- & 8/8 & E & $<$0.001 & significant \\
C$'$-matched $-$ B-matched & 5 & +2.72 & +2.86, +3.08, +2.50, +2.94, +2.22 & 5/5 & S & 0.004 & significant (exploratory family) \\
D $-$ R1 & 5 & +78.77 & --- & 5/5 & S & $<$0.001 & significant (exploratory family) \\
E1-last $-$ D, lifetime & 8 & $-$24.02 & $-$25.37, $-$25.15, $-$24.53, $-$23.21, $-$24.06, $-$22.67, $-$21.79, $-$25.35 & 0/8 (E1 ahead) & S & unadj.\ $3\times10^{-10}$ & exploratory; not pre-registered \\
E2 $-$ D, lifetime & 8 & $-$47.02 & $-$49.56, $-$48.59, $-$46.79, $-$48.21, $-$46.29, $-$44.25, $-$42.46, $-$50.01 & 0/8 (E2 ahead) & S & unadj.\ $3\times10^{-10}$ & exploratory; not pre-registered \\
\bottomrule
\end{tabular}
\end{table*}

\subsection{Tuning and compute explain most of the first headline}
\label{sec:tuning}
Our first draft reported D 83.3 against C$'$ 49.3 and B 38.6.
The burst learning rate was at the edge of the calibration grid for B and C$'$.
At the tuned rate, B rises from 38.59 to 83.32 (+44.73, 3 of 3 seeds, Holm 0.003).
C$'$ rises from 49.30 to 84.31 (+35.01), and D from 83.33 to 88.52 (+5.19).
The lead of D over C$'$ shrinks from 34.03 points (old rate, unmatched), to 4.21 (tuned, unmatched), to 0.97 (tuned and step-matched, 3 seeds) and 1.06 (5 seeds).
On eight seeds it is $-0.005$.
Compute matching and tuning changed the conclusion at every step.
Figure~\ref{fig:tuning} shows the arms before and after tuning.
The same pattern holds on VGG6 and ViT-tiny, where tuning gains for B are +44.56 and +24.17 points.

\begin{figure}[t]
\centering
\figorbox{figs/fig7_tuning.pdf}{\columnwidth}{2.1in}
\caption{Untuned versus tuned baselines (ResNet-9, CIFAR-10, n=3 seeds: 42, 101, 102). Each point is one seed. The burst learning rate moves from 0.005 (old) to the tuned value. B rises from 38.59 to 83.32 and C$'$ from 49.30 to 84.31 on the mean; D rises from 83.33 to 88.52. Notice that the controller's lead over C$'$ shrinks from 34.03 to 4.21 as the baselines are tuned.}
\label{fig:tuning}
\end{figure}

We picked the ResNet-9 rates on seed-42 \emph{test} accuracy, on the same 5,000 images we report.
That choice is optimistic.
Table~\ref{tab:lr} also gives the choice from the canary set.
The largest loss for any arm is 0.45 points, and D's rate does not change.
The test pick helped the baselines, not D.
D has more settings than the matched rivals, which have only a learning rate: a gate margin, a plateau rule and priors (taken from one seed's checkpoint).
This favors D, so the lack of an edge over the rivals is not due to an under-tuned D.
VGG6 and ViT-tiny rates were picked on canary accuracy only.
Matched arms exist only at rate 0.5, so the canary-picked matched gaps are projections: D minus B-matched shifts by at most +0.45, and D minus C$'$-matched by +0.05.

\begin{table}[t]
\centering
\caption{Learning-rate selection, ResNet-9, 3 seeds. Both picks are shown. Loss is the 3-seed accuracy lost by using the canary-picked rate instead of the test-picked rate.}
\label{tab:lr}
\small
\begin{tabular}{lccc}
\toprule
Arm & Test pick & Canary pick & Loss (pt) \\
\midrule
B & 0.5 & 0.3 & 0.45 \\
C$'$ & 0.5 & 0.3 & 0.06 \\
D & 0.3 & 0.3 & 0.00 \\
\bottomrule
\end{tabular}
\end{table}

\subsection{The matched rival: pre-registered replication}
\label{sec:replication}
The strongest fair rival is C$'$-matched.
It has living spares and D's burst steps.
On the first five seeds, D was ahead of it by 1.06 points (95\% interval +0.5 to +1.6, 5 of 5 seeds).
The Holm-adjusted $p$ was 0.056, so we labelled the result consistent, not significant.
We also did not quote the sign test, because seeds 103 and 104 were added after our first look.
That interim result is the reason we planned E3.

E3 is a pre-registered replication on three new seeds, 105 to 107.
The rules were fixed before the runs (commit 018276a).
\emph{Success} needs a positive mean, at least 7 of 8 seeds ahead, a Holm $p$ below 0.05, and all three new seeds ahead.
\emph{Failure} is a mean below +0.5 or at most 5 of 8 seeds ahead.
The outcome is a failure.
The per-seed gaps on the new seeds are $-3.12$, $-3.32$ and $+1.10$ points (mean $-1.78$; D ahead on 1 of 3).
Over all eight seeds the mean gap is $-0.005$, D is ahead on 6 of 8, and the Holm-adjusted $p$ in the recomputed headline family is 1.000.
Figure~\ref{fig:replication} shows all eight seeds.

We draw three conclusions.
First, the +1.06 lead is not a claim of this paper, and we do not call it replicated.
The five seeds that showed it were the seeds that prompted E3, so the estimate was likely too high.
Second, the defensible headline against matched rivals is no detectable difference on 8 seeds, not a proven tie.
We ran no equivalence test.
The mean of $-0.005$ mixes two groups: D is ahead on 6 of 8 seeds by 0.44 to 1.52 points, and behind by 3.12 and 3.32 points on the 2 seeds where the gate skipped a burst.
Third, the E3 test is a sequential design with an interim look: three seeds were added after a five-seed analysis.
Any $p$ value from the eight-seed set is optimistic, so we print the new seeds alone and the pooled set.
The fix for the gate (a fixed pre-deployment reference) is untested, so we can neither claim that D works when fixed nor dismiss it.
Recomputing the Holm family with the eight-seed contrast leaves the D versus B-matched result significant ($p=0.0014$).

\begin{figure}[t]
\centering
\figorbox{figs/fig6_replication.pdf}{\columnwidth}{2.1in}
\caption{Per-seed difference D minus C$'$-matched in final accuracy, all eight seeds. The five original seeds (42, 101--104) gave +1.28, +0.44, +1.18, +1.52, +0.88 (mean +1.06). The three pre-registered new seeds (105, 106, 107; E3) gave $-3.12$, $-3.32$, +1.10 (mean $-1.78$). The eight-seed mean is $-0.005$, with D ahead on 6 of 8 seeds. The numbers under each seed are D's burst steps (1475, 1500, 1375, 1500, 1500, 775, 700, 1500). Notice that the two strongly negative new seeds are the low-step ones.}
\label{fig:replication}
\end{figure}

\subsection{Reviewer baselines: E1 and E2}
\label{sec:e1e2}
Reviewers will ask two questions.
Does the closest prior art, adapted to waves, match D?
Does a forward-only BatchNorm repair already recover most of the gain?
E1 and E2 answer them.
Their rules were registered before the runs.

\emph{E1 is not distinguishable from D on final accuracy.}
E1 runs one Liu-style burst at the last wave, with a greedy saliency-ordered remap each wave.
Its burst has D's step count and a rate of 0.5 from an earlier plateau; we did not re-sweep it.
E1-last minus D is $-0.55$ points over eight seeds (D ahead on 6 of 8; E1 ahead on seeds 105 and 106), Holm $p=0.66$.
The pre-registered label is mixed: D and E1 cannot be told apart on 8 seeds, and we ran no equivalence test.
The result is heterogeneous: D is ahead on 6 seeds by 0.88 to 2.06 points, and E1 is ahead by 2.02 and 2.82 points on seeds 105 and 106, where D's gate skipped a burst.
E1 is far above remap-only (+76.18, 8 of 8 seeds, significant), so the sanity check passes.
A single burst at the first wave (E1-first) collapses once later waves arrive: finals are 10.08 to 27.88 percent, and the mean is 68.11 points below D (0 of 8 seeds).
A burst before the last wave cannot cover later faults.
On final accuracy, a simple one-shot retrain at D's compute is not distinguishable from D's paced, gated repair on this test.
It does not say that E1 is the better design: E1 has no repair between waves, and it ran with a rate that was not swept.

\emph{E1 does not serve through life.}
Final accuracy is the registered measure for E1, but a deployed chip must serve between waves.
We therefore computed lifetime accuracy (mean of the three per-wave accuracies) for E1-last and E2 on the same runs.
This analysis is exploratory, not pre-registered, and in family S.
E1-last minus D is $-24.02$ points over 8 seeds (D ahead on 8 of 8).
E2 minus D is $-47.02$ points (D ahead on 8 of 8).
The mechanism is plain.
E1 repairs wave 0 by greedy remap, then the pool is empty, and with no retraining before the last wave it serves at 10 to 17 percent after wave 1.
Its single last-wave burst brings it back to 83.3 to 87.7 percent.
This gap is close to by construction for a one-shot method, and we do not read it as D beating Liu's method.
C$'$-matched also retrains in service, and it is not distinguishable from D on lifetime ($-0.17$, mixed; Section~\ref{sec:spares}).
The lesson is that one terminal retrain matches D on the final number but not on service through life.

\emph{E2 stays at chance when the spare pool is exhausted.}
E2 remaps at random and refits BatchNorm.
It takes no optimizer step.
Wave 0 is repaired (89.12 and 88.90 percent on seeds 103 and 104).
Then the spares run out and accuracy falls.
The eight finals range from 8.58 to 13.64 percent, with a mean of 10.58.
E2 minus D is $-76.06$ points (0 of 8 seeds, negative).
E2 minus R1 is +0.67 (5 of 8 seeds, Holm $p=0.66$, mixed): BatchNorm statistics alone do not repair the collapse.
The pre-registered rule says that if E2 came within 2 points of D, we must withdraw the claim that bursts are needed.
E2 is far outside that range, so the claim stands: burst retraining is required in this scenario.
E2 ran at 27.1\% cumulative faults and 15\% spares, where remap-only is also at chance because the pool runs out.
We therefore claim only that BatchNorm refit does not rescue a chip whose spare pool is exhausted (the tested regime).
We did not test it at a low fault rate.

\subsection{The gate has a drift failure mode}
\label{sec:gate}
D lost to C$'$-matched on seeds 105 and 106 (D 81.24 and 80.76 against 84.36 and 84.08).
We read the per-wave logs after the verdicts were fixed.
This diagnosis is post-hoc; it changes no verdict, no seed and no gate definition.
The canary gate measures the drop against the canary reading after the previous wave's repair.
On both seeds the wave-1 burst stopped early at a plateau (275 and 200 steps), and left a weak reference (83.69 and 83.40 percent).
At wave 2 the canary dropped by only 1.76 and 0.20 points, below the gate threshold of 2.0, so D skipped the burst and ended near 81 percent.
On a typical seed (42) the drop at that wave was 7.62 points and D ran all three bursts (500, 500 and 475 steps).
The reference moves down after a weak burst, so the bar moves down with it.
We read this as a reference-drift failure mode, seen in 2 of 8 seeds, and report it as a limitation.
C$'$-matched (fixed burst length, no gate) and E1 (one fixed burst) are not exposed to it.
D's under-spending also gave the step-matched rivals fewer steps on those seeds, yet they finished 2.0 to 3.3 points higher.
The causal reading is a judgment from one trace per seed, not a registered test.
A gate that measures against the pre-deployment canary would be a new design, and it needs its own pre-registered test.
We did not run it, so the gate fix is untested.

\subsection{The overwrite effect}
\label{sec:overwrite}
At the old burst rate (0.005), living spares gained 2.48 points over frozen references (46.82 against 49.30 percent; 3 of 3 seeds; Figure~\ref{fig:overwrite}).
Both arms ended at 47 to 49 percent there, so that rate barely worked, and we report the number only as an old-rate result that is superseded.
We then ran a pre-registered control at the tuned rate (0.5; 3 seeds; on a different GPU, with fault maps identical across arms within a seed).
Living minus frozen final accuracy is $+0.03$ points (per seed $+0.08$, $0.00$, $0.00$), which is noise-level, and the pre-registered decision was NOT SUPPORTED.
A third arm that starts the spare from its original value also lies within noise.
The lifetime mean accuracy is $0.31$ points lower for living spares, on 3 of 3 seeds.
Frozen ran 0.42 to 0.96 points above living after the first wave, and the last burst brought the arms together.
The mechanism (a remap restores the original reference over a retrained cell) is a code-level observation, and we did not show that it costs accuracy at the tuned rate.
The design also cannot test the case it is about.
In the capacity policy, every remap happens before the first retraining burst, so the value at fault of each remapped cell equals its original programmed value.
The C versus C$'$ contrast (old and new rate) therefore compares only a pinned spare with a trainable one; it never tests overwriting learned compensation.
Only D remaps after retraining, and it has no isolated ablation of that step.
This result fits the paper's main thesis: a tuned, matched comparison erased an apparent gain, here the living-spares gain, as it erased D's lead.
The stuck-weight freeze has no isolated measurement for quantized devices.

\begin{figure}[t]
\centering
\figorbox{figs/fig2_overwrite.pdf}{\columnwidth}{2.1in}
\caption{The overwrite bug. Per-wave accuracy for frozen-reference spares (C) and living spares (C$'$) at the old burst rate (n=3 seeds: 42, 101, 102). Final accuracy is 46.82 for C and 49.30 for C$'$ (+2.48; per seed +2.38, +2.64, +2.42; paired $t$ test $p=0.0011$). This old-rate effect is superseded: at the tuned rate the difference is $+0.03$ points (noise-level; Section~\ref{sec:overwrite}).}
\label{fig:overwrite}
\end{figure}

\subsection{Spares and lifetime}
\label{sec:spares}
Remap alone is a capacity cliff.
At 2, 5, 10 and 15\% spares it ends at 9.9, 11.8, 10.0 and 10.3 percent; at 30\% it equals the fault-free chip at 89.02.
The cliff sits at the cumulative stuck fraction of 27.1\%, by construction.
We did not run spare counts between 15 and 30\%.
On the three seeds with a 30\% R1 run, D at 15\% spares ends at 88.62 percent, within 0.4 points of the clean chip (89.01; remap-only at 30\% is 89.02; interval $-1.0$ to $+0.2$, 0 of 3 seeds ahead).
On all eight seeds D ends at 86.64 percent.
D uses about 1,475 burst steps to do this.
This is a fixed-schedule result at one network and one schedule, and it trades compute for area.
We do not read it as a spare-area saving.
Figure~\ref{fig:cliff} shows the sweep.
Against C$'$-matched, D's gap is noise-level at 2\% spares (+0.49), mixed at 5\% (+0.70), and consistent at 10\% (+2.33), all with 3 seeds.

\begin{figure}[t]
\centering
\figorbox{figs/fig5_spare_cliff.pdf}{\columnwidth}{2.1in}
\caption{Spare sweep. Remap only (R1, 3 seeds) is at chance for 2, 5, 10 and 15\% spares (9.9, 11.8, 10.0, 10.3 percent) and equals the fault-free chip at 30\% (89.02); the cliff is the cumulative stuck fraction of 27.1\%, by construction. D at 15\% spares (n=3) ends at 88.62 percent. Notice that D spends about 1,475 burst steps to do so. The 88.62 value is a 3-seed mean; on eight seeds D ends at 86.64.}
\label{fig:cliff}
\end{figure}

Pacing has a price over the lifetime.
Figure~\ref{fig:lifetime} shows accuracy after each wave.
On the five original seeds, D serves 86.09, 87.44 and 88.50 percent after waves 1, 2 and 3.
C$'$-matched serves 88.92, 86.18 and 87.44.
D's lead over C$'$-matched on those seeds exists only at end of life, and it did not replicate (Section~\ref{sec:replication}).
Averaged over the lifetime, D minus C$'$-matched is $-0.17$ points (2 of 5 seeds, mixed).
D minus B-matched is +4.32 points (5 of 5 seeds, Holm 0.0027, significant).
Pacing holds spares back, so D misses the strict clean-minus-2 floor after wave 1 in all five seeds.
In an ablation with equal steps, pacing adds about 0.7 points (3 seeds, consistent, confounded: the ablated arm fires two fixed bursts and D fires three adaptive ones).
After the E3 result we do not read this as an accuracy gain of pacing or gating over step-matched baselines.
Pacing is a lifetime trade-off: it lowers early accuracy and does not raise final accuracy over a matched rival.

\begin{figure}[t]
\centering
\figorbox{figs/fig4_lifetime.pdf}{\columnwidth}{2.1in}
\caption{Accuracy after each wave (mean over seeds; band is the range over seeds). D and C$'$-matched: eight seeds (D 86.02, 86.48, 86.64; C$'$-matched 88.97, 85.62, 86.64). B-matched: five seeds (80.73, 83.61, 84.72). R1: eight seeds (88.97, 12.71, 9.90). Notice the trade-off: pacing costs D accuracy after wave 1, and after wave 3 the D and C$'$-matched means are equal. The five-seed values in the text differ from these eight-seed means.}
\label{fig:lifetime}
\end{figure}


\subsection{Several designs, the same final accuracy: what each costs}
\label{sec:tradeoffs}
Table~\ref{tab:tradeoffs} puts the designs that reach about the same final accuracy side by side.
Matching D on final accuracy is itself a result: it tells a reader that the choice can be made on cost, not on a margin we could not show.
The rows differ in what they need and in how they fail, not in a ranking.

\begin{table*}[t]
\centering
\caption{Designs that reach the same final accuracy at matched compute, and what each costs (ResNet-9, CIFAR-10, three waves of 10\%, 15\% spares). Final accuracy is the mean over $n$ seeds, in percent. Lifetime is the mean over the three per-wave accuracies, given as the arm minus D in points; the E1-last row is exploratory and not pre-registered. ``Settings'' lists the choices a user must make, as our implementation exposes them; it is a qualitative list, not a measured quantity. This is a comparison of costs and failure modes, not a ranking.}
\label{tab:tradeoffs}
\footnotesize
\begin{tabular}{@{}>{\raggedright\arraybackslash}p{1.55cm}>{\raggedright\arraybackslash}p{0.95cm}>{\raggedright\arraybackslash}p{1.45cm}>{\raggedright\arraybackslash}p{0.9cm}>{\raggedright\arraybackslash}p{1.2cm}>{\raggedright\arraybackslash}p{2.4cm}>{\raggedright\arraybackslash}p{2.2cm}>{\raggedright\arraybackslash}p{4.3cm}@{}}
\toprule
Design & Final acc. ($n$) & Lifetime vs D & Spares? & Canary / held-out set? & Settings to tune & When it retrains & Known failure mode \\
\midrule
D (controller) & 86.64 (8) & --- & Yes, living & Yes (gate and plateau stop) & Several: burst rate, gate threshold, plateau rule, expected waves, BN batches & After a wave, in an adaptive burst, if the canary drops & Gate reference drifted and skipped a needed burst on 2 of 8 seeds (about 3 points behind on those) \\
C$'$-matched (capacity trigger) & 86.64 (8) & +0.17 ($n=5$, mixed) & Yes, living & No & Burst rate and burst length (set to D's step budget) & In a fixed burst when a wave overflows the spare pool & None observed in the main test; living spares gave +0.03 over frozen at the tuned rate (noise-level) \\
E1-last (one Liu-style retrain) & 86.08 (8) & $-24.02$ ($n=8$) & Yes & No & Burst rate and length; which wave & Once, at the last wave & Serves at 10--17\% between waves once the pool is empty; D ahead on 6 of 8 seeds, not distinguishable \\
B-matched (retrain only) & 84.72 (5) & $-4.32$ ($n=5$) & No & No & Burst rate and length & After every wave, fixed burst & 3.78 points below D (5 of 5 seeds); not re-run on seeds 105--107 \\
R1 (remap only) & 9.90 (8) & Collapses after wave 1 & Yes & No & None & Never & At chance once the pool is exhausted (E2, a BatchNorm refit, does too: 10.58) \\
\bottomrule
\end{tabular}
\end{table*}

Which design to choose depends on what the deployment can afford.
If only final accuracy matters, the simplest design that matches D is C$'$-matched: it needs no canary set and no gate, and it is not exposed to the gate drift we found.
One terminal retrain (E1-last) is enough only if accuracy in service between waves does not matter, because it serves at 10--17 percent until the last wave.
If the lifetime curve matters, retraining in service is needed, and C$'$-matched and D are not distinguishable on it.
Spares together with in-service retraining are the necessary core: retraining alone ends 3.78 points lower, and remap alone ends at chance once the pool is exhausted.
What D adds is a canary-driven trigger and paced spare use, and at matched compute we did not measure that they buy accuracy; the gate also has the drift failure mode above.

\subsection{Other networks, fault patterns and stress cases}
\label{sec:stress}
Table~\ref{tab:other} lists the other settings.
These rows are separate three- or five-seed runs.
They are not a general edge over C$'$-matched, and the E3 result applies to all of them.
On CIFAR-100, D's gap over the best C$'$-matched rate is +7.16 points (3 of 3 seeds, indicative); the rates were picked on test accuracy of seed 42, and the C$'$ rate sweep may overstate the gap.
On realistic fault models (Weibull timing, stuck-at plus drift, line faults), the gaps are +0.94, +1.11 and +0.73 points (3 of 3 seeds each, consistent).
On ViT-tiny the gap is mixed (+1.50, 4 of 5 seeds).
On VGG6 it is mixed ($-0.52$, 3 of 5 seeds): the edge does not reproduce.

D has no edge or trails in three settings.
With many small waves and a gate fixed to the best-so-far canary, D trails C$'$-matched at equal steps: $-2.22$ (six waves) and $-0.41$ (eight waves), both mixed.
On hardware-calibrated scenarios (S4, small steady drift), D trails B-matched by 0.65 to 3.54 points; the two cases with 5\% spares are negative (0 of 3 seeds).
We believe, but have not tested, that the canary gate and plateau stop end bursts too early under small steady drift.
The seed 105 and 106 traces point to the same family of failure (Section~\ref{sec:gate}).
The drift size and timing in S4 are our assumptions and drive the result.
The D2 variant (an adaptive trigger that stops bursts earlier), written down before any run, addressed this.
It passed its calibrated-scenario criterion in 4 of 4 scenarios and failed its main-test criterion, so D stays the controller.

\begin{table*}[t]
\centering
\caption{D against C$'$-matched and B-matched in other settings (final accuracy, points). Seeds: 3 unless stated. Status follows Section~\ref{sec:methodology}. Rates for CIFAR-100 were picked on test accuracy of seed 42. These runs are separate from the eight-seed main test and are not a general margin over C$'$-matched.}
\label{tab:other}
\small
\begin{tabular}{p{4.6cm}p{0.8cm}p{3.1cm}p{3.1cm}p{3.6cm}}
\toprule
Setting & $n$ & D $-$ C$'$-matched & D $-$ B-matched & Status (vs C$'$-matched) \\
\midrule
2\% spares & 3 & +0.49 (3/3) & +0.78 & noise-level \\
5\% spares & 3 & +0.70 (2/3) & +2.37 & mixed \\
10\% spares & 3 & +2.33 (3/3) & +3.41 & consistent \\
25\% faults per wave & 3 & +2.14 (3/3) & +3.74 & consistent \\
CIFAR-100 & 3 & +7.16 (3/3) & --- & consistent (df=2) \\
VGG6 & 5 & $-$0.52 (3/5) & +2.58 (4/5) & mixed \\
ViT-tiny & 5 & +1.50 (4/5) & +8.20 (5/5) & mixed (vs C$'$-matched) \\
Weibull timing (S1) & 3 & +0.94 (3/3) & +3.05 & consistent \\
Stuck-at plus drift (S2) & 3 & +1.11 (3/3) & +3.51 & consistent \\
Line faults (S3) & 3 & +0.73 (3/3) & +2.79 & consistent \\
Six small waves, gate fixed & 3 & $-$2.22 (1/3) & --- & mixed \\
Calibrated drift S4b, 15\% spares & 3 & --- & $-$1.05 (1/3) & mixed \\
Calibrated drift S4b, 5\% spares & 3 & $-$4.95 (0/3) & $-$3.54 (0/3) & negative \\
\bottomrule
\end{tabular}
\end{table*}

\subsection{What the parts contribute}
We ablated D's parts (3 seeds).
Importance ranking beats random order by 1.65 to 1.73 points, consistent.
Ranking is prior art and a design choice here.
Fisher saliency is no better than weight magnitude (+0.00 and +0.29, mixed).
Dropping the BatchNorm refit changes the result by $-0.08$ points (mixed, old rate only).
Removing the gate and plateau stop costs 0.79 points at the old rate; we never ran this at the tuned rate.
We therefore claim none of these parts as a gain.
The part that carries weight in our evidence is the retraining bursts (Section~\ref{sec:e1e2}); living spares showed no tuned-rate accuracy effect (Section~\ref{sec:overwrite}).

\subsection{Closed loop on a probe-estimated fault map}
\label{sec:closedloop}
All main results use an oracle fault map.
We replaced it with a map estimated from counted, noisy probe reads.
The estimator and its registered tests belong to the companion paper (in preparation).
Here we use it as a black box.
It reads each 256 by 256 tile with random $\pm 1$ probes, decodes a sparse set of stuck cells, and falls back to a full readout when decoding hits a cap.
We ran D on the estimated map and on the oracle map with the same seed and the same fault realization (42 pairs, equal fault-map hashes).
Registered bars were non-inferiority at $-1.0$ points (dense) and $-0.5$ points (sparse).

At 2\% read noise, we define the gap as D on the estimated map minus D on the oracle map.
The mean gap is $-0.04$ points in a dense regime and $-0.07$ in a sparse regime (5 seeds each; ranges $-0.42$ to $+0.24$ and $-0.24$ to $+0.02$).
Both are above the registered bars of $-1.0$ and $-0.5$ points.
On the 3-seed registered run, the largest per-seed gap is 0.24 points.
One identical oracle rerun differed by 0.24 points, and rerun noise elsewhere has an RMS of 0.46 points, so we cannot separate these gaps from noise.
This is not an equivalence test.

Three limits apply.
First, the dense regime is a full-readout fallback in 99\% of tiles, at 1.12 times the reads of a full readout.
It shows that a measured map suffices, not that probing localizes.
Only the sparse regime tests probing.
Second, the registered reads bar (at most 0.25 of a full readout on every seed) held on the 3-seed run.
On five seeds, one seed used 0.2525, so that clause was not met.
Third, at 4\% noise the sparse extension was killed: the worst-wave recall of new stuck cells was 0.853 to 0.894 against a bar of 0.90.
Accuracy still held (+0.06 points sparse, +0.30 dense).
The result is a no-difference-detected result with 3 to 5 seeds, not an equivalence test.
It covers stuck-at faults, one tile size and 2\% to 4\% read noise.
It addresses the oracle-map objection only within that scope.

% =================================================================================================
\section{Discussion and Limitations}
\label{sec:discussion}

\subsection{What the evidence supports}
The evidence supports four statements.
First, a remap that restores a frozen reference can erase learned compensation.
At the tuned rate, living spares showed no measurable accuracy effect (+0.03 points, 3 seeds, pre-registered, NOT SUPPORTED); the earlier +2.48 was an old-rate artefact, and our design cannot test remap after retraining.
Second, retraining is needed once the pool is exhausted.
A forward-only BatchNorm refit stays at chance on 8 of 8 seeds, 76.06 points below D, in the tested regime only.
Third, at equal burst compute and a tuned rate, D beats retraining-only by 3.78 points, significant after Holm correction on five seeds (a sequential design: seeds were extended from 3 to 5 after a first look).
Fourth, D is not distinguishable from the best matched living-spares rival (eight seeds, mean $-0.005$) or from a one-shot Liu-style retrain at D's compute ($-0.55$, Holm $p=0.66$) on final accuracy; we ran no equivalence test, and the result is heterogeneous (D ahead on 6 of 8 seeds, behind by about 3 points on 2).
The one-shot retrain serves far below D through life (exploratory, $-24.02$ points over 8 seeds).
These statements hold only in simulation.

The most useful result for other groups may be the evaluation record.
A 34-point lead fell to 4.2 points when we tuned the rival's rate, to 1.06 when we matched compute, and to $-0.005$ when a pre-registered replication added three seeds.
The living-spares control points the same way: a gain of 2.48 points at an old rate vanished at the tuned rate.
A repair controller needs a tuned, compute-matched baseline before its margin means anything.
It also needs a replication, because our own five-seed result of +1.06 points did not survive one.

\subsection{Why simple matched baselines are not distinguishable from the controller}
We did not test why D is not distinguishable from them.
Three readings fit the data, and each is a hypothesis.
First, once spares are living and the burst has D's compute, the order and the pace of the remap may matter little.
Second, a canary gate that reads a moving reference can skip a burst that a fixed schedule would run (Section~\ref{sec:gate}).
Third, C$'$-matched and E1 get the same step count as D with no way to under-spend.
Our evidence does not rank these readings.
Practitioners can take one lesson: spend the retraining compute. Keeping learned weights when a spare is remapped is sensible by construction, but we did not show that it matters for accuracy.

\subsection{Limitations}
\begin{itemize}
\item \textbf{Simulation only.}
Every result is from one simulator (CrossSim) with stuck-at faults injected by us.
There is no device measurement.
Main runs have no read noise, programming error or drift.
The total stuck rate and the HRS to LRS mix in S4 come from one second-hand report; timing is assumed.
\item \textbf{No controller margin over matched rivals.}
D has no measured margin over C$'$-matched (eight seeds, mean $-0.005$, Holm $p=1.0$) or over E1-last (Holm $p=0.66$) on final accuracy.
We ran no equivalence test, and the results are heterogeneous (D ahead on 6 of 8 seeds, behind by about 3 points on the 2 gate-skip seeds).
The five-seed +1.06 is an interim-look result that did not replicate.
A single-seed gap under 1.22 points is within rerun noise.
\item \textbf{Retraining-only is a five-seed result.}
B-matched was not run on seeds 105 to 107, so the +3.78 gap rests on five seeds, extended from three after a first look.
D's two weak seeds (105, 106) would lower it; we did not measure how much.
\item \textbf{The gate has a drift failure mode.}
The diagnosis is post-hoc and rests on two seeds.
A repaired gate (a fixed pre-deployment reference) was not tested.
Canary-gated retraining and the plateau stop have prior art and no measured benefit at the tuned rate.
\item \textbf{Sequential design.}
Seeds 103 and 104 were added after a first look, and E3 added seeds 105 to 107 after a five-seed analysis.
\item \textbf{The network matters.}
D's edge over C$'$-matched is mixed on VGG6 and ViT-tiny.
CIFAR-100 and ViT-tiny have three and five seeds.
\item \textbf{Calibrated drift is negative.}
D trails retraining-only in S4 (negative with 5\% spares, mixed with 15\%), and the cause is an untested hypothesis.
\item \textbf{Oracle map in the main results.}
Closed loop at 2\% read noise narrows this; the 4\% extension was killed.
\item \textbf{Controller inputs.}
D needs the expected number of waves.
Its priors came from one seed's checkpoint, and one priors file moved a single-seed result by 7.7 points.
\item \textbf{Cost.}
Retraining compute and write endurance are not costed; all matched arms spend the same compute.
Labelled data must be available on the device or at the edge.
\item \textbf{Baselines.}
E1's burst rate was not re-swept.
An earlier, simpler Liu-style reimplementation (remap, then plain retraining, rate swept) trailed D by 2.03 points on three seeds; it is not the faithful comparison and is not used as evidence. Its 0.3--0.5 rate plateau is only the source of E1's fixed rate.
Retraining-free compensation schemes were not run.
\item \textbf{Selection and provenance.}
ResNet-9 rates were picked on test data (both picks printed).
D has more settings than the matched rivals, which have only a learning rate.
Pooled seeds mix fault maps from two GPU models.
No run records its code commit, and one snapshot commit holds the code used.
\end{itemize}

\subsection{Next steps}
A pre-registered test of a gate that measures against the pre-deployment canary would show whether the drift failure mode costs accuracy.
We would also run B-matched on seeds 105 to 107 under a registered plan, sweep E1's burst rate, test the overwrite in a policy that remaps after retraining (the tuned-rate control is done), and run spare counts between 15 and 30\%.

% =================================================================================================
\section{Conclusion}
\label{sec:conclusion}
A chip that loses cells over its life spends a finite spare budget.
We found, by code reading, that remapping to a frozen reference restores the original value over what retraining had learned.
Living spares gave no measurable accuracy gain at the tuned rate (+0.03 points, 3 seeds, pre-registered; an earlier +2.48 at an old rate is superseded), and our design could not test remap after retraining.
This agrees with the main finding: tuned, compute-matched baselines erase apparent gains.
A forward-only BatchNorm refit does not rescue a chip whose spare pool is exhausted (the tested regime).
At equal compute and tuned rates, our controller D beats retraining-only by 3.78 points, in 5 of 5 seeds, significant after Holm correction; the seeds were extended from 3 to 5 after a first look.
D is not distinguishable from the matched living-spares rival or a one-shot Liu-style retrain on final accuracy (8 seeds, no equivalence test).
The result is heterogeneous: D is ahead on 6 of 8 seeds and behind by about 3 points on the 2 seeds where its gate skipped a burst, and the gate fix is untested.
The one-shot retrain serves far below D through life (exploratory, 8 seeds), so one terminal retrain matches D on the final number but not on service through life.
At matched compute we found no accuracy gain from pacing or gating.
The 34-point lead of our first draft shrank to zero under compute matching, tuning and a pre-registered replication.
With a probe-estimated map at 2\% read noise, D stayed within the registered non-inferiority bars of its oracle-map accuracy (not an equivalence test), and D has no edge or trails on calibrated drift and with many small waves.
The condensed catalogue of simulator errors may help other groups avoid the same ones.

% =================================================================================================
\section*{Availability}
Code, configurations, raw run records and analysis scripts: \url{https://github.com/prateek1656/NeuroFault}. Commit identifiers cited in this paper refer to the author's timestamped development repository, available to reviewers on request.

% =================================================================================================
\bibliographystyle{IEEEtran}
\bibliography{refs}

\end{document}
